\documentclass[11pt]{article}

\usepackage[french]{babel}

\usepackage[T1]{fontenc}

\usepackage{amsfonts}
\usepackage{amsmath}
\usepackage{amssymb}

\usepackage{graphicx}
\usepackage{multirow}
\usepackage{multicol}
\usepackage{xcolor}
\usepackage[shortlabels]{enumitem}

\usepackage[margin=2cm]{geometry}

\usepackage{mathtools}
\usepackage{siunitx}
\usepackage{physics}

\usepackage{tikz}
\usepackage{hyperref}

\usepackage{algorithm}
\usepackage{algpseudocode}
\usetikzlibrary{quantikz2,babel}

\tikzset{qcircuit/.style={
    column sep=.9cm,
    row sep=0.7cm,
    operator/.style={draw=black!80, line width=0.5pt, fill=gray!5, minimum height=0.7cm, minimum width=0.7cm},
    phase/.style={draw=black!80, fill=black!20, circle, minimum size=4pt},
    slice/.style={font=\footnotesize\itshape, color=gray!50, dashed}
}}


\DeclarePairedDelimiter\ceil{\lceil}{\rceil}
\DeclarePairedDelimiter\floor{\lfloor}{\rfloor}


\AtBeginDocument{\RenewCommandCopy\qty\SI}

\makeatletter
   \def\vhrulefill#1{\leavevmode\leaders\hrule\@height#1\hfill \kern\z@}
\makeatother

\newcommand{\moy}[1]{\langle #1 \rangle}
\newcommand{\pDer}[2]{\frac{\partial #1}{\partial #2}}
\newcommand{\pDerd}[2]{\frac{\partial^2 #1}{\partial #2^2}}

\newcounter{questioncounter}
\newenvironment{question}[1]{
\refstepcounter{questioncounter}
\noindent\rule{0.05\textwidth}{1pt}\,\,\raisebox{-.3\ht\strutbox}{\textbf{Question \thequestioncounter \hspace{5mm} #1}}\,\,\vhrulefill{1pt}\par\vspace{0.4cm}
}{\par}

%======================================================
%en-tete
%======================================================
\begin{document}
\selectlanguage{french}
\noindent{\Large Anthony Wroblewski} \hfill {\Large \today}\\
{\Large Nathan Brochu}

\par\vspace{5mm}
\begin{center}
   {\Large PHQ344 - Physique Statistique 1} \\ \vspace{5mm}
   {\LARGE\sffamily Devoir \#2} \\ \vspace{5mm}
\end{center}

%======================================================
%question 1
%======================================================
\begin{question}{Ensemble microcanonique}
Soit un système microcanonique de $N = 4$ spins avec des moments magnétiques $\mu_i = \pm \mu$ et une énergie totale $E = +2\mu H$.
\\\\
L'énergie totale est donnée par $E = -H \sum_{i=1}^4 \mu_i = +2\mu H$, ce qui implique que la somme des moments est :
\begin{equation*}
\sum_{i=1}^4 \mu_i = -2\mu
\end{equation*}
Pour obtenir une somme égale à $-2\mu$ avec 4 spins de valeur $\pm\mu$, il faut nécessairement avoir 3 spins avec un moment $-\mu$ et 1 spin avec un moment $+\mu$.
\\\\
Le nombre total de microétats accessibles possédant cette énergie est le nombre de façons de placer un spin $+\mu$ parmi 4 emplacements :
\[
\Omega = \binom{4}{1} = 4 \quad \text{donc} \quad \Omega = 4
\]
Les microétats correspondants sont $(-1, -1, -1, +1)\mu$, $(-1, -1, +1, -1)\mu$, $(-1, +1, -1, -1)\mu$ et $(+1, -1, -1, -1)\mu$.
\\\\
La probabilité qu'un spin individuel ait un moment $-\mu$ est :

\begin{equation*}
P(-\mu) = \dfrac{3}{4}
\end{equation*}
\\
La probabilité qu'un spin individuel ait un moment $+\mu$ est : 
\begin{equation*}
P(+\mu) = \dfrac{1}{4}
\end{equation*}


L'aimantation moyenne du système se calcule par la moyenne de la somme des moments magnétiques sur tous les microétats accessibles :
\[
\text{Aimantation moyenne} = \frac{(-1 - 1 - 1 + 1)\mu}{4} = -\frac{1}{2}\mu
\]
\end{question}
\newpage
%======================================================
%question 2
%======================================================
\begin{question}{Thermodynamique d'un gaz imparfait}
On recherche la variation d'entropie $\Delta S = S_1 - S_0$. D'après le premier principe de la thermodynamique :
\[
\delta E = \delta Q + \delta W
\]
où $\delta Q$ représente le transfert thermique et $\delta W$ le travail transféré.
\bigskip

Pour une transformation réversible :
\begin{itemize}[label=$$]
    \item $\delta Q = T \, dS$ (car la température $T$ est constante)
    \item $\delta W = -F \, dx = -P \, dV $
\end{itemize}

\bigskip 
En substituant ces expressions dans le premier principe :

\begin{equation*}
\delta E = T \, dS - P \, dV
\end{equation*}
De l'équation $PV = N k_B (T + b T^2)$, on tire la pression $P = \dfrac{N k_B (T + b T^2)}{V}$. Ainsi :
\[
\delta E = T \, dS - \frac{dV}{V} N k_B (T + b T^2)
\]
En isolant $dS$ :
\[
dS = \frac{1}{T} \, \delta E + \frac{dV}{V} N k_B (1 + b T)
\]
En intégrant entre les états initial et final :
\[
\Delta S = \frac{1}{T} \int \delta E + N k_B (1 + b T) \int_{V_0}^{V_1} \frac{dV}{V}
\]
\[
\Delta S = \frac{1}{T} \int \delta E + N k_B (1 + b T) \ln\left(\frac{V_1}{V_0}\right)
\]

On sait par l'énoncé que l'énergie $E$ est indépendante du volume $V$ et de la pression $P$. Comme la température $T$ est également fixe ($T = \text{cte}$), l'énergie reste constante au cours du processus, soit $dE = 0$.
\[
\delta E = \delta Q + \delta W \implies \delta E = T \, dS - P \, dV = 0
\]
Par conséquent, le terme d'intégration de l'énergie s'annule ($\int \delta E = 0$), ce qui nous donne directement :
\[
\Delta S = N k_B (1 + b T) \ln\left(\frac{V_1}{V_0}\right)
\]

\end{question}
\newpage
%======================================================
%question 3
%======================================================
\begin{question}{Réfrigérateur à absorption}
\subsubsection*{a) Contrainte (Conservation de l'énergie)}
Considérons le cycle de fonctionnement de la machine ($\Delta E = dQ = 0$). Les chaleurs entrantes sont $Q_f$ (flamme) et $Q_r$ (réfrigérateur), tandis que la chaleur rejetée vers l'atmosphère est $Q_a$.
\[
Q_a = Q_f + Q_r
\]

\subsubsection*{b) Coefficient de performance $\eta_{\max}$}
Le second principe de la thermodynamique impose que la variation totale d'entropie lors d'un cycle soit supérieure ou égale à zéro : $\Delta S_{\text{tot}} \ge 0$. L'entropie des réservoirs varie selon :
\[
\Delta S = \frac{Q_a}{T_a} - \frac{Q_r}{T_r} - \frac{Q_f}{T_f} \ge 0
\]
En réarrangeant les termes :
\[
\frac{Q_a}{T_a} \ge \frac{Q_r}{T_r} + \frac{Q_f}{T_f}
\]
En utilisant la relation $Q_a = Q_r + Q_f$:
\[
\frac{Q_r + Q_f}{T_a} \ge \frac{Q_r}{T_r} + \frac{Q_f}{T_f}
\]
En divisant toute l'inégalité par $Q_f$ :
\[
\frac{Q_r}{Q_f} \frac{1}{T_a} + \frac{1}{T_a} \ge \frac{Q_r}{Q_f} \frac{1}{T_r} + \frac{1}{T_f}
\]
Puisque $\eta = \dfrac{Q_r}{Q_f}$, on exprime :
\[
\frac{1}{T_a} - \frac{1}{T_f} \ge \eta \left(\frac{1}{T_r} - \frac{1}{T_a}\right)
\]
On isole $\eta$ :
\[
\eta \le \left(\frac{1}{T_a} - \frac{1}{T_f}\right) \left(\frac{1}{T_r} - \frac{1}{T_a}\right)^{-1}
\]
La valeur maximale du coefficient de performance est donc :
\[
\eta_{\max} = \left(\frac{1}{T_a} - \frac{1}{T_f}\right) \left(\frac{1}{T_r} - \frac{1}{T_a}\right)^{-1}
\]

\subsubsection*{c) $T_f$ élevée}
En réécrivant la formule de $\eta_{\max}$ en mettant en valeur le facteur contenant $T_f$ :
\[
\eta_{\max} = \left(\frac{1}{T_a} - \frac{1}{T_f}\right) \left(\frac{T_r T_a}{T_a - T_r}\right)
\]
Pour des températures $T_a$ et $T_r$ fixes avec $T_a > T_r$, le terme $\left(\frac{T_r T_a}{T_a - T_r}\right)$ est une constante positive.

Pour maximiser $\eta_{\max}$, on veut que le terme $\left(\frac{1}{T_a} - \frac{1}{T_f}\right)$ soit le plus grand possible. Cela nécessite de rendre le terme $\dfrac{1}{T_f}$ le plus petit possible, c'est-à-dire d'avoir $T_f$ la plus grande possible ($\frac{1}{T_f} \to 0$).

\end{question}
\newpage
%======================================================
%question 4
%======================================================
\begin{question}{Questions courtes}

\subsubsection*{a)}
D'après le second principe de la thermodynamique ($\Delta S \ge 0$), la chaleur circule spontanément du corps le plus chaud vers le corps le plus froid.
\[
\Delta S = -\frac{Q_{\text{ext}}}{T_{\text{ext}}} + \frac{Q_{\text{fluide}}}{T_{\text{fluide}}} \ge 0
\]
Selon le premier principe ($\Delta E = Q_{\text{ext}} - Q_{\text{fluide}} = 0$), on a $Q_{\text{ext}} = Q_{\text{fluide}}$. Ainsi :
\[
\frac{Q_{\text{ext}}}{T_{\text{fluide}}} \ge \frac{Q_{\text{ext}}}{T_{\text{ext}}} \implies T_{\text{ext}} \ge T_{\text{fluide}}
\]
Pour absorber de la chaleur de l'environnement extérieur ($Q_{\text{ext}} > 0$), le fluide doit être plus froid que l'air extérieur ($T_{\text{fluide}} < T_{\text{ext}}$) afin que le transfert thermique s'effectue spontanément vers le fluide.

\subsubsection*{b)}
\begin{itemize}\itemsep0.1em
    \item Chaleur absorbée : $Q_{100} = 2.0\,\text{J}$ à $T_{100} = 100\,^\circ\text{C} = 373.15\,\text{K}$
    \item Chaleur rejetée : $Q_0 = 1.0\,\text{J}$ à $T_0 = 0\,^\circ\text{C} = 273.15\,\text{K}$
    \item Travail produit : $W = 1.0\,\text{J}$
\end{itemize}
\bigskip

Il faut que la conservation de l'énergie soit respectée :
\[
\Delta E = Q_{100} - Q_0 - W = 2.0 - 1.0 - 1.0 = 0\,\text{J}
\]
Le premier principe est respecté ($Q_{100} - Q_0 = W$).
\bigskip

Il faut aussi que l'entropie augmente :
\[
\Delta S = -\frac{Q_{100}}{T_{100}} + \frac{Q_0}{T_0} \ge 0 \implies \frac{Q_0}{T_0} \ge \frac{Q_{100}}{T_{100}}
\]
Étant donné que $Q_0 = 1.0\,\text{J}$ et $Q_{100} = 2.0\,\text{J}$, on a $Q_0 = \frac{1}{2} Q_{100}$. En substituant :
\[
\frac{\frac{1}{2} Q_{100}}{T_0} \ge \frac{Q_{100}}{T_{100}} \implies T_{100} \ge 2 T_0
\]
Vérifions numériquement :
\[
373.15\,\text{K} \ge 2 \cdot 273.15\,\text{K} \implies 373.15\,\text{K} \ge 546.3\,\text{K}
\]
Puisque $373.15\,\text{K} \ge 546.3\,\text{K}$ est une inégalité fausse, la condition du second principe $\Delta S \ge 0$ n'est pas respectée. Cette machine n'est pas possible selon les lois de la thermodynamique car elle ne satisfait pas le second principe.
\end{question}


\end{document}